\begin{thebibliography}{99}
\bibitem[1]{B} S. Billey, Kostant polynomials and the cohomology ring for G/B, Duke Math. J. 96 (1999), 205–224.
\bibitem[2]{BM} A. S. Buch and L C. Mihalcea, Curve neighborhoods of Schubert varieties, J. Differential Geom. 99 (2015), no. 2, 255–283.
\bibitem[3]{FRT} J. S. Frame, G. de B. Robinson, and R. W. Thrall, The hook graphs of the symmetric group, Can. J. Math. 6 (1954), 316–325.
\bibitem[4]{G} E. R. Gansner, The Hillman–Grassl correspondence and the enumeration of reverse plane partitions, J. Combin. Theory Ser. A 30 (1981), 71–89.
\bibitem[5]{GK} W. Graham and V. Kreiman, Excited Young diagrams, equivariant K-theory, and Schubert varieties, Trans. Amer. Math. Soc. 367 (2015), 6597–6645.
\bibitem[6]{H} J. E. Humphreys, Reflection Groups and Coxeter Groups, Cambridge Stud. Adv. Math., vol. 29, Cambridge Univ. Press, 1992.
\bibitem[7]{IN1} T. Ikeda and H. Naruse, Excited Young diagrams and equivariant Schubert calculus, Trans. Amer. Math. Soc. 361 (2009), 5193–5221.
\bibitem[8]{IN2} , K-theoretic analogues of factorial Schur P- and Q-functions, Adv. Math. 243 (2013), 22–66.
\bibitem[9]{IT1} M. Ishikawa and H. Tagawa, Schur function identities and hook length posets, Proceedings of the 19th International Conference on Formal Power Series and Algebraic Combinatorics (Tianjin, July 2–6, 2007), 2007, available at http://igm.univ-mlv.fr/\{\textasciitilde{}\}fpsac/FPSAC07/SITE07/ PDF-Proceedings/Posters/55.pdf.
\bibitem[10]{IT2} , Leaf posets and multivariate hook length property, RIMS Kokyuroku 1913 (2014), 67– 80.
\bibitem[11]{Ka} M. Kashiwara, The flag manifold of Kac–Moody Lie algebra, in Algebraic Analysis, Geometry, and Number Theory: Proceedings of the JAMI Inaugural Conference (J. Igusa, ed.), Johns Hopkins Univ. Press, 1989, pp. 161–190.
\bibitem[12]{KY} J. Kim and M. Yoo, Hook length property of d-complete posets via q-integrals, J. Combin. Theory, Ser. A 162 (2019), 167–221.
\bibitem[13]{KN} A. N. Kirillov and H. Naruse, Construction of double Grothendieck polynomials of classical types using IdCoxeter algebras, Tokyo J. Math. 39 (2017), no. 3, 695–728.
\bibitem[14]{Knu} D. E. Knuth, The Art of Computer Programming, Volume 3: Sorting and Searching, 3rd Edition, Addison-Wesley, 1973.
\bibitem[15]{Kr1} V. Kreiman, Schubert classes in the equivariant K-theory and equivariant cohomology of the Grassmannian, https://arxiv.org/abs/math/0512204.
\bibitem[16]{Kr2} , Schubert classes in the equivariant K-theory and equivariant cohomology of the Lagrangian Grassmannian, https://arxiv.org/abs/math/0602245.
\bibitem[17]{Ku} S. Kumar, Kac–Moody Groups, their Flag Varieties and Representation Theory, Prog. Math., vol. 204, Birkhäuser, 2002.
\bibitem[18]{LSS} T. Lam, A. Schilling, and M. Shimozono, K-theory Schubert calculus of the affine Grassmannian, Comp. Math. 146 (2010), no. 4, 811–852.
\bibitem[19]{LP} C. Lenart and A. Postnikov, Affine Weyl groups in K-theory and representation theory, Int. Math. Res. Not. IMRN 2007 (2007), rnm038.
\bibitem[20]{LS} C. Lenart and M. Shimozono, Equivariant K-Chevalley rules for Kac–Moody flag manifolds, Amer. J. Math. 136 (2014), no. 5, 1175–1213.
\bibitem[21]{M} L. Mihalcea, On equivariant quantum cohomology of homogeneous spaces: Chevalley formulae and algorithms, Duke Math. J. 140 (2007), no. 2, 321–350.
\bibitem[22]{MPP} A. Morales, I. Pak, and G. Panova, Hook formulas for skew shapes I. q-analogues and bijections, J. Combin. Theory Ser. A 154 (2018), 350–405.
\bibitem[23]{N3} K. Nakada, q-Hook formula for a generalized Young diagram, in preparation.
\bibitem[24]{N1} , Colored hook formula for a generalized Young diagram, Osaka J. Math. 45 (2008), no. 4, 1085–1120.
\bibitem[25]{N2} , q-Hook formula of Gansner type for a generalized Young diagram, 21st International Conference on Formal Power Series and Algebraic Combinatorics (FPSAC 2009) (Hagenberg, Austria), DMTCS Proceedings, vol. AK, Discrete Math. Theor. Comput. Sci., 2009, pp. 685–696.
\bibitem[26]{Naruse} H. Naruse, Schubert calculus and hook formula, Talk slides at 73rd Sém. Lothar. Combin., Strobl, Austria, 2014, available at https://www.mat.univie.ac.at/\textasciitilde{}slc/wpapers/s73vortrag/ naruse.pdf.
\bibitem[27]{Okada} S. Okada, (q, t)-Deformations of multivariate hook product formulae, J. Algebraic Combin. 32 (2010), no. 3, 399–416.
\bibitem[28]{PY} O. Pechenik and A. Yong, Equivariant K-theory of Grassmannians, Forum Math. $\Pi$ 5 (2017), e3.
\bibitem[29]{Proc1} R. A. Proctor, Dynkin diagram classification of $\lambda$-minuscule Bruhat lattices and d-complete posets, J. Algebraic Combin. 9 (1999), no. 1, 61–94.
\bibitem[30]{Proc2} , Minuscule elements of Weyl groups, the number game, and d-complete posets, J. Algebra 213 (1999), no. 1, 272–303.
\bibitem[31]{Proc3} , d-Complete posets generalize Young diagrams for the hook product formula: Partial presentation of proof, RIMS Kokyuroku 1913 (2014), 120–140.
\bibitem[32]{PS} R. A. Proctor and L. M. Scoppetta, d-Complete posets: Local structural axioms, properties, and equivalent definitions, Order (2018), https://doi.org/10.1007/s11083-018-9473-4.
\bibitem[33]{Schur} I. Schur, Über die Darstellung der symmetrischen und der alternierenden Gruppe durch gebrochene lineare Substitutionen, J. Reine Angew. Math. 139 (1911), 155–250.
\bibitem[34]{Stan2} R. P. Stanley, Theory and application of plane partitions, Part 2, Studies in Applied Math. 50 (1971), no. 3, 259–279.
\bibitem[35]{Stan1} , Ordered Structures and Partitions, Mem. Amer. Math. Soc., vol. 119, Amer. Math. Soc., 1972.
\bibitem[36]{Stan3} , Enumerative Combinatorics, Volume I, Cambridge Stud. Adv. Math., vol. 49, Cambridge Univ. Press, 1997.
\bibitem[37]{Stem1} J. R. Stembridge, On the fully commutative elements of Coxeter groups, J. Algebraic Combin. 5 (1996), no. 4, 353–385.
\bibitem[38]{Stem2} , Minuscule elements of Weyl groups, J. Algebra 235 (2001), no. 2, 722–743.
\bibitem[39]{T} R. W. Thrall, A combinatorial problem, Michigan Math. J. 1 (1952), no. 1, 81–88.
\end{thebibliography}
